\section{Introdução}
\label{sec:introducao}

A seção de introdução é a porta de entrada do seu artigo. Nela, você deve apresentar o contexto geral em que a pesquisa se insere, preparando o leitor para o problema que será abordado. Esta seção deve convencer o leitor sobre a relevância do tema. Ao longo do texto, utilize citações de outros artigos para embasar suas afirmações (exemplo de citação de artigo: \cite{artigo_exemplo}).

\subsection{Descrição do Problema}
Nesta subseção, você deve detalhar qual é exatamente o problema que o seu trabalho visa resolver ou investigar. O problema deve ser formulado de maneira clara e objetiva. Por que esse problema existe e quais as suas consequências?

\subsection{Justificativa}
A justificativa deve responder à pergunta: "Por que este trabalho é importante?". Aqui, você descreve a relevância científica, tecnológica e social da sua pesquisa, destacando os benefícios que a resolução do problema trará para a área de estudo ou para a sociedade.

\subsection{Objetivo principal e específicos}
O \textbf{objetivo principal} define, em uma frase ampla, o que você pretende alcançar com este artigo. 
Os \textbf{objetivos específicos} são os passos intermediários ou metas menores necessárias para que o objetivo principal seja alcançado. Geralmente são apresentados em formato de tópicos:
\begin{itemize}
    \item Realizar uma revisão sistemática da literatura;
    \item Implementar o protótipo XYZ;
    \item Validar o modelo através de experimentos.
\end{itemize}

\subsection{Público Alvo}
\textit{Nota: Em artigos científicos tradicionais, o "público-alvo" geralmente fica implícito na justificativa ou na introdução. Contudo,  esta subseção pode ser incluída para especificar quem se beneficiará diretamente com a solução proposta (pesquisadores, desenvolvedores, empresas do setor X, etc).}

\subsection{Organização das seções do artigo}
O restante deste artigo está organizado da seguinte forma. A Seção~\ref{sec:fundamentacao} apresenta os conceitos teóricos necessários. A Seção~\ref{sec:trabalhos_relacionados} discute trabalhos correlatos. A Seção~\ref{sec:metodologia} detalha a metodologia adotada. Na Seção~\ref{sec:analise_resultados}, são apresentados e discutidos os resultados experimentais. Por fim, a Seção~\ref{sec:consideracoes_finais} conclui o trabalho e aponta direções futuras.
